\section{Fixed central conditioning and sufficient source cost}
\label{sec:source-conditioning}
The construction above is finite, but its expanding core floor can force
$\log(1/\varepsilon)$ to grow quadratically with the joining radius.
Note~55 \cite{n55} removes that sufficient loss from the growing part
of the core. The positive regulator is inverted only on a fixed central
rectangle; direct source collars handle the rest. This changes the
sufficient source description bound while preserving the physical
regulator, field prescription and complete tuple.

\subsection{A fixed central region and a direct source comparison}
First fix $b$, the physical parameter $\eta$ and the scalar reserves.
For the all-time construction $b=T_++1$ is already selected. Choose
\[
 R_0=2\max(R_b^*,2e^{\Lambda_b}+32)+64,\quad
 e_0=\min(1,E_be^{-(R_0+16)^2/72}),\quad R_c\ge2R_0+64.
\]
Use the compact positive-regulator transfer on $|x|\le R_0+8$ with
coefficients and mean caps evaluated at $R_0+16$. If $D_r^0$ are
its finite row coefficients, including auxiliary mean rows, let
\begin{equation}
 K_0=1+\frac4{e_0^2}\max_r D_r^0.
 \label{eq:sc-central}
\end{equation}
Its source error is at most $K_0\varepsilon$ per row. This argument
allows multiple zeros. The coefficient can be extremely large, but it
is independent of the growing $R_c$.

Put $R=R_c+8$ and
\begin{equation}
 \Psi_b(R)=\pi R/8+A_b^{\rm amp}[1+\log(2+R)]^2,
 \qquad \zeta=\varepsilon e^{\Psi_b(R)+32}.
 \label{eq:sc-carrier-parameter}
\end{equation}
The armwise full-buffer bound \eqref{eq:at-carrier} pays this relative
source parameter throughout the remaining core and bands. For the
sharper fixed $[95,97]$ table one can instead use
\begin{equation}
 \Psi_{97}(R)=\pi R/8+12\log^2(2+R)+200\log(2+R)+202.
 \label{eq:sc-carrier97}
\end{equation}
The added $32$ covers fixed carrier/collar comparisons.
Assume $\zeta\le e^{-1}$ and $\zeta\log(2+R)\le2^{-512}$.

For $\Delta=G-H$ compare $G_s=H+s\Delta$ directly with $H$ on
its proved outer collars. Rouch\'e on a larger disk and then a small
root circle, followed by unit comparison, gives the conservative table
\begin{center}\small
\begin{tabular}{@{}ll@{}}\toprule
Quantity & Bound, with $\ell_c=\log(2+|a_H|)$\\\midrule
Root displacement & $|a_s-a_H|\le2^{64}\zeta/\ell_c$\\
Normalized slope and unit & between $2^{-64}$ and $2^{64}$\\
Log unit ratio; its derivative & $2^{96}\zeta$; $2^{112}\zeta\ell_c$\\
Residue difference; individual residues & $2^{160}\zeta\ell_c$; $32$\\
Analytic remainder difference in $G_s''/G_s$ & $2^{192}\zeta\ell_c^2$\\\bottomrule
\end{tabular}\end{center}
For the residue use jets through three at the paired roots; for the
regular part subtract both poles on the larger boundary before using
analyticity. $P$ has supplied the geometry of $H$, but its field error
is not an additive term in $G_s-H$. This is essential: the source column
must tend to zero as $\varepsilon\to0$ at a fixed model and seam.

With $\ell_R=\log(2+R)$, the arbitrary-positive-$q$ radial estimates give
\begin{equation}
 s[1],w[1]\le2^{1024}\zeta,\quad
 v_m[1]\le2^{1024}\zeta[1+\log(1/\zeta)+\ell_R^2],\quad
 z_m[1]=0\quad\hbox{in the pure exact-heat core}.
 \label{eq:sc-primitives}
\end{equation}
To see the logarithm, remove scaled disks of radius $O(\zeta)$ about
the two roots. Their pole integral is $O(\zeta)$. On the complement
root/unit interpolation bounds the pole-kernel variation by
$C\zeta(|u-a|^{-2}+|u-a|^{-1})$; radial integration gives
$\zeta\log(1/\zeta)$. Local parameters can be smaller than $\zeta$;
use monotonicity of $u[1+\log(1/u)]$ on $(0,1]$, not an upper bound
on their logarithms alone.
Every collar meeting the split at $R_0$ is included whole, with its
fixed overlap paid centrally. Only nonnegative raw integrals are split;
they are recombined before weak differentiation. No new interface is
introduced into the actual field.

\subsection{Band variation and the complete consumer}
On spatial bands use the proof family
$F_s=G_s+\lambda(P-G_s)$. Its source directions are exactly
\begin{align}
 \delta F&=(1-\lambda)\Delta,&
 \delta B&=-\lambda_x\Delta/2,\notag\\
 \delta E&=-\lambda_x\Delta_z-\lambda_{xx}\Delta/4-\lambda_t\Delta,
 \notag\\
 \delta E_z&=-\lambda_x\Delta_{zz}
 -3\lambda_{xx}\Delta_z/4-\lambda_{xxx}\Delta/8
 -\lambda_t\Delta_z-\lambda_{xt}\Delta/2.
 \label{eq:sc-band-directions}
\end{align}
Time directions of $B,B_z$ use $\Delta_t=-\Delta_{zz}$ and
$\Delta_{zt}=-\Delta_{zzz}$. Thus four spatial jets suffice even for
the moving gap. Although the component $\lambda E_P$ has zero source
direction, its regularized $Z_m$ varies with the actual field and weight;
that contribution is retained. Only the pure exterior source column is
identically zero.

On the normalized gap floor, $u=\overline F/(|F|^2+q)$ satisfies
\[
 \delta u=\frac{\overline{\delta F}}{|F|^2+q}
 -u\frac{2\Re(\delta F\overline F)}{|F|^2+q},\quad
 |\delta u|\le3|\delta F|/|F|^2,\quad
 |\delta w|\le2w(1-w)|\delta F/F|.
\]
Insert these identities in the exact compatibility, heat and stress
formulas. The variations of $T_C,\Xi$ cost $C\zeta\ell_R^2$;
those of current and potential corrections cost
$C\zeta(\ell_R^3+\ell_R^2)$; $\delta(T_C)_t$ costs
$C\zeta[\ell_R^4+|r_t|\ell_R^2]$.
These are pointwise band estimates before area $O(1/\ell_R)$.
No inverse physical regulator enters them.

The physical coefficient caps at height $u$ have degrees at most two;
the largest composed weight, $|r_t|R_1^2$, has degree four. It is used
only on source-active supports, where $u\le R$.
First pay the auxiliary mean rows. Actual means then have the fixed
central cap or $C_{H,b}(1+\ell_R)+1$ on annular and seam supports.
Apply the signed-stress identity before absolute bounds and the actual
global mean-product subtraction afterwards. This controls $Q,\Delta Q$,
not covariances assembled independently from field pieces.

For an explicit common coefficient let $W_J$ dominate the physical
coefficients after their stated height powers are removed, $L_{\rm prof}$
the fixed profile constants after removal of $\ell_R^j$, and $D_{\rm mod}$
the residual coefficient. The choices
\[
 W_J=2^{16}(1+A_b+B_b'),\quad
 L_{\rm prof}=2^{64}(1+e^{4/3}),\quad D_{\rm mod}=2^{32}(1+b)^4
\]
suffice. With the complete moments and mean coefficient already fixed, set
\begin{align}
 A_J={}&1+K_0+C_{H,b}+\sum_{j=0}^3(B_j(b)+1)+W_J+L_{\rm prof}
                        +e_0^{-1}+R_0+K_*+D_{\rm mod},\notag\\
 C_J^{\rm src}={}&2^{8192}A_J^{128},\qquad P_R=(1+R)^8.
 \label{eq:sc-coefficient}
\end{align}
The fixed-central coefficients have degree less than sixteen in these
inputs; band variation has denominator powers at most four after
normalization and at most three profile derivatives. The finite transfer
maps are linear except for the mean product. Even thirty-two factors
and $2^{16}$ terms fit \eqref{eq:sc-coefficient}. The degree-eight height
envelope covers the degree-four physical weight and remaining profile,
pole and mean logarithms. Thus every source-active scalar row obeys
\begin{equation}
 \mathcal E_{{\rm source},r}
 \le K_0\varepsilon+\tfrac12 C_J^{\rm src}P_R
                    \zeta[1+\log(1/\zeta)]
 \le C_J^{\rm src}P_R\zeta[1+\log(1/\zeta)].
 \label{eq:sc-source-map}
\end{equation}
This includes $U,W$ through four, the tuple and direct time/source rows
through two, internal current order three, the heat-potential correction
through four, and the paid traces. The source-free exterior and the
nonsource seam/model reserves remain additional. On the late temporal
core one uses the separate relative field premise
\eqref{eq:ep-source-premise}, rather than treating its blend as exact heat.

For $0<\theta_s\le1$ a finite inverse of the modulus is
\begin{equation}
 d_s=\frac{\theta_s}{C_J^{\rm src}P_R},\quad
 H_s=1+\log(4/d_s),\quad
 \zeta_s=\frac{d_s}{4H_s},\quad
 \varepsilon_{\rm loc}=e^{-\Psi_b(R)-32}\zeta_s.
 \label{eq:sc-tolerance}
\end{equation}
Indeed $1+\log(1/\zeta_s)=H_s+\log H_s\le2H_s$, so the
source allowance is at most $\theta_s/2$; the large fixed coefficient
also pays the root/unit gates. Choose
$\theta_s=\min(\delta/4,B_b^+(X_0))$ for the all-time family.
Then $-\log\theta_s=O_b(\log R)$ as $R\to\infty$. A fixed nonzero
source reserve would not prove a vanishing total spatial error.

\subsection{Replace every growing source constraint}
On the fixed late central rectangle define
\[
 A_*=e^{-4b\log^2(b+2)-2(R_0+5)},\qquad c_*=\min(1,A_*).
\]
The full late core has carrier lower bound $c_*e^{-\Psi_b(R)}$:
centrally use $A_*$ and outside use \eqref{eq:at-carrier}.
At the spatial gap $e^{-\Psi_b(R)}$ alone suffices. Hence one can
replace \emph{all three} coarse growing demands in
\eqref{eq:at-source-minimum} by
\begin{equation}
 \varepsilon_{\rm new}=\min\left\{1,\varepsilon_{\rm loc},
 a_X^{-2}e^{-\Psi_b(R)},T_-^{-5/2},
 \frac{\delta}{4(1+K_{\rm mid}(T_-))},
 \frac{\delta}{2^{69}K_{\rm br}L^3},
 \frac{e_{\min}c_*e^{-\Psi_b(R)}}{16}\right\}.
 \label{eq:sc-all-source}
\end{equation}
Keeping an old expanding core floor or coarse late carrier clause in
this minimum would retain its larger exponent. The last entry pays the
late field premise separately; observable transfer alone cannot pay it.
All field formulas, temporal compatibility terms, actual mean caps and
both endpoint comparisons of Theorem~\ref{thm:all-time} are retained.
The original source rule remains valid and may be preferable at a
particular finite radius.

\subsection{An explicit source size and its asymptotic meaning}
For any fixed upper time $s\ge0$ and $0<\varepsilon\le1$ put
$q_s=\log(1/\varepsilon)$, $Q_s=\log(32/\varepsilon)$ and
\begin{align}
 U&=\lceil4+\log(2+s+Q_s)\rceil,& F_s&=U^4e^{sU^2+4U},\notag\\
 n_0&=\max\left(3,\left\lceil\sqrt{\log(2^{14}F_s/\varepsilon)/2}\right\rceil\right),
 &N&=n_0-1,\notag\\
 M&=\left\lceil U\max\{1,(2^{20}F_s/(48\varepsilon))^{1/4}\}\right\rceil,
 &h_s&=U/M,\notag\\
 \beta_{\rm bits}&=\lceil\log_2(64F_s/\varepsilon)\rceil,
 &B_{\rm bits}&=\beta_{\rm bits}+20+\lceil\log_2(M+1)\rceil.
 \label{eq:sc-constructor}
\end{align}
The same dyadic Hermite values/slopes and positive clipping give at most
$4M$ pieces and the four-jet error $29\varepsilon/64$, for all $t\le s$,
all real heights and strip four. For the shorter tail cutoff,
$s,Q_s\le e^{U-4}$ and $e^U\ge2U^2$ imply
\[
 sU^2+13U-\pi e^{4U}\le-2e^{4U},\qquad
 4\pi e^{4U}-2sU-13-j/U\ge\pi e^{4U}\quad(j\le3).
\]
The full tail is therefore at most $\varepsilon/8$; the interpolation,
omitted-mode and nodal-arithmetic debits are unchanged. Clipping is a
contraction toward the positive true source.

An explicit descriptor upper bound is
\begin{equation}
 S_{\rm bound}=64(M+1)(B_{\rm bits}+1)\quad\hbox{bits}.
 \label{eq:sc-storage}
\end{equation}
It stores rational data and defining polynomials/isolation labels for
algebraic clipping endpoints. Since the maximum in the mesh formula is
its second entry,
\[
 M=\left\lceil(2^{20}/48)^{1/4}U^2
                   e^{q_s/4+sU^2/4+U}\right\rceil.
\]
For fixed $s$ and $q_s\to\infty$ this proves
\begin{align}
 \log M,\ \log S_{\rm bound}
   &=q_s/4+(s/4)\log^2 q_s+O_s(\log q_s),\notag\\
 \beta_{\rm bits},B_{\rm bits}&=O_s(q_s),\qquad N=\Theta_s(\sqrt{q_s}).
 \label{eq:sc-q-cost}
\end{align}
Choose the displayed tolerance, rather than arbitrarily over-refining it.
Equations~\eqref{eq:sc-tolerance}--\eqref{eq:sc-all-source} give, with
$b,\eta,\theta$ fixed,
\begin{equation}
 q_s=\pi R/8+O_{b,\eta,\theta}(\log^2R),\qquad
 \log M,\ \log S_{\rm bound}
       =\pi R/32+O_{b,\eta,\theta}(\log^2R).
 \label{eq:sc-radius-cost}
\end{equation}
For the sharper $[95,97]$ policy the log-square coefficient in the last
expression is $(12+97)/4=109/4$; for the general fixed-$b$ policy it is
$(A_b^{\rm amp}+b)/4$. The earlier expanding-floor policies had leading
log cell bounds $R_c^2/904$ on $[95,97]$ and $R_c^2/144$ for the
coarse $[63,b]$ table. These compare specified sufficient constructions,
not necessary complexity. When $b$ changes with accuracy, use the finite
formulas \eqref{eq:sc-constructor}--\eqref{eq:sc-storage}; the fixed-$b$
asymptotics do not control that joint limit.

The node algorithm in Note~55 has a rational arithmetic-term bound
$O_s(Mq_s^{11/2})$. Clipping can use a pruned Sturm tree: a degree-three
cell has at most three occupied intervals per depth, and coefficient
height at most $2^{B_{\rm bits}}$ gives sufficient isolation depth
$8+4B_{\rm bits}$. Repeated cubic roots are rational and are handled
separately by the fixed-degree gcd. Thus a conservative extra bound is
$2^{20}M(B_{\rm bits}+1)$ rational operations, without scanning the
whole fine grid. This is an operation count, not bit complexity.

Each piece contributes Gaussian moments through degree six to $G$ and
its first three derivatives, so the symbolic prescription uses $O(M)$
fixed-order primitive calls per point. For this clipped-source alternative, accuracy of complex error
functions and algebraic endpoints needs a separate implementation ledger.
Sections~\ref{sec:analytic-source} and~\ref{sec:complete-evaluation}
provide a compressed positive polynomial alternative with paid source
cancellation, observable arithmetic and continuous-filter error/cost
analysis. They do not presume stability of this predecessor's primitives.
The source-data bits in \eqref{eq:sc-q-cost} do not bound the working
precision of that evaluator. No source coefficients or new numerical
experiments have been evaluated for this integration.
